\documentclass[11pt]{article}
\usepackage[margin=0.82in]{geometry}
\usepackage[T1]{fontenc}
\usepackage{lmodern,amsmath,amssymb,amsthm,microtype}
\usepackage[colorlinks=true,urlcolor=blue]{hyperref}
\newtheorem{proposition}{Proposition}
\newcommand{\Q}{\mathbb Q}
\newcommand{\Z}{\mathbb Z}
\title{Denominator bounds and rational telescopers\\for symmetric zeta forms}
\author{}
\date{}
\begin{document}
\maketitle
\begin{abstract}
We derive a uniform denominator bound for the normalized symmetric kernels,
construct their rational telescoping certificates after lower-period
cancellation, and retain the full gcd in the resulting arithmetic cost.
The denominator theorem is a paper proof. Lean verifies the leading
binomial identity, block recurrence, positivity implication, and invariance
under excess clearing. No all-odd irrationality theorem is claimed.
\end{abstract}

\section{An explicit denominator budget}
Fix an odd $s=2r+1\ge3$, an even $n\ge0$, and $0\le m\le n$. Define
\[
 T_{n,m}(t)=\frac{(n!)^s}{(m!)^2}
 \frac{\prod_{v=1}^m(t-v)(t+n+v)}
      {\prod_{j=0}^n(t+j)^s},\qquad
 d_n=\operatorname{lcm}(1,\ldots,n),\quad d_0=1.
\]
This proper rational function has a unique expansion
\[
 T_{n,m}(t)=\sum_{j=0}^n\sum_{q=1}^s\frac{c_{j,q}}{(t+j)^q}.
\]
\begin{proposition}
For every displayed $j,q$,
\begin{equation}
 d_n^{s-q}c_{j,q}\in\Z.
 \label{eq:budget}
\end{equation}
Consequently the convergent sum has the form
\[
 \mathcal A_{n,m}:=\sum_{k\ge1}T_{n,m}(k)
   =b_{n,m}+\sum_{\substack{3\le q\le s\\q\ {\rm odd}}}
                  a_{q,n,m}\zeta(q),
\]
where $d_n^{s-q}a_{q,n,m}\in\Z$ and $d_n^s b_{n,m}\in\Z$.
\end{proposition}
\begin{proof}
Factor $T_{n,m}=P_m Q_n^s$, where
\[
 P_m(t)=\binom{t-1}{m}\binom{t+n+m}{m},
 \qquad Q_n(t)=\frac{n!}{\prod_{h=0}^n(t+h)}.
\]
For any integer $a$, Vandermonde's polynomial identity and
\[
 \binom{a+x}{m}=\sum_{i=0}^m\binom{a}{m-i}\binom{x}{i},
 \qquad
 \binom{x}{i}=\frac{(-1)^{i-1}x}{i}
                  \prod_{h=1}^{i-1}(1-x/h)\quad(i\ge1)
\]
show that $d_m^\ell$ clears its coefficient of $x^\ell$. Indeed each
contributing denominator is a product of $\ell$ integers from $1,\ldots,i$.
The coefficients $\binom a k$ are integral also for negative $a$, by
$\binom{-b}{k}=(-1)^k\binom{b+k-1}{k}$. Products retain the weighted
Taylor bound. Since $m\le n$, it holds with $d_n$ for $P_m(-j+x)$.

At a pole,
\[
 [xQ_n(-j+x)]^s=(-1)^{js}\binom nj^s
                    \prod_{h\ne j}(1+x/(h-j))^{-s}.
\]
The coefficient of $x^\ell$ in each factor is
$(-1)^\ell\binom{s+\ell-1}{\ell}(h-j)^{-\ell}$; $d_n^\ell$ clears it.
Thus the coefficient of $x^{s-q}$ in
$P_m(-j+x)[xQ_n(-j+x)]^s$ satisfies \eqref{eq:budget}. That coefficient
is exactly $c_{j,q}$.

Reflection gives $c_{n-j,q}=(-1)^{q+1}c_{j,q}$, canceling even zeta
coefficients. Proper decay gives $\sum_jc_{j,1}=0$, and partial-fraction
summation yields
\[
 a_{q,n,m}=\sum_jc_{j,q},\qquad
 b_{n,m}=-\sum_{j=0}^n\sum_{q=1}^s c_{j,q}H_j^{(q)}.
\]
For $q=1$, subtract $H_K\sum_jc_{j,1}=0$ from the $K$-th partial sum;
the remaining upper boundary tends to zero since
$0\le H_{K+j}-H_K\le n/(K+1)$. No value of $\zeta(1)$ is used.
Finally $d_n^q H_j^{(q)}\in\Z$, proving the assertion about $b_{n,m}$.
The case $n=0$ follows directly from $T_{0,0}(t)=t^{-s}$.
\end{proof}

This is an explicit application to our kernels of the
integral-valued-polynomial arithmetic appearing in \cite[Section 2]{Zudilin}.
It reduces the clearing range from $n+m$ to $n$. In particular,
\[
 c_{j,s}=(-1)^{m+j}\binom nj^s\binom{j+m}{m}\binom{n-j+m}{m}.
\]
Lean verifies the factorial-to-binomial magnitude identity and integrality.
Its identification with the Laurent coefficient is part of the paper proof.

\section{The combined form and its actual arithmetic cost}
Use the $r$ blocks $m=n-j$, $0\le j<r$, for sufficiently large even $n$.
Choose nonzero integer weights $w_j$ canceling the $r-1$ lower odd periods.
Put $M_n=\sum_j|w_j|$ and
\[
 J_n=d_n^s\sum_jw_j\mathcal A_{n,n-j}=A_n\zeta(s)+B_n,\qquad
 h_n=\gcd(|A_n|,|B_n|).
\]
When $h_n>0$, the primitive form is $J_n/h_n$, with the bound
\begin{equation}
 \left|\frac{J_n}{h_n}\right|
 \le K_n\max_j\mathcal A_{n,n-j},
 \qquad K_n=\frac{d_n^s M_n}{h_n}.
 \label{eq:cost}
\end{equation}
This keeps the cancellation weights and the complete gcd in the same
calculation. Replacing the clearing factor by $c\,d_n^s$ multiplies
$A_n,B_n,h_n$ by $c$, so $K_n$ is unchanged. Lean verifies this invariance:
sharper clearing can improve an upper bound that ignores the gcd, but
cannot alter an already primitive form.

The companion note's saddle theorem \cite[Section 5]{Companion} gives,
for every fixed $j$,
\[
 \frac1n\log\mathcal A_{n,n-j}\longrightarrow\lambda_s<0
 \quad(n\longrightarrow\infty,\ n\ {\rm even}).
\]
For $s=5$, $\lambda_5\in(-4.970794,-4.970793)$. Hence a sufficient
arithmetic estimate would be
\[
 \limsup_{n\to\infty,\ n\ {\rm even}}\frac{\log K_n}{n}<-\lambda_s.
\]
This strict inequality is not established here. Nonvanishing is also
required. At positive integers $T_{n,m}(k)=0$ for $k\le m$ and
$T_{n,m}(k)>0$ thereafter. Therefore weights of one sign give a nonzero
combined sum. Lean proves the corresponding convergent-series implication;
the existence of one-sign nullvectors at all indices is not proved.

\section{A rational telescoping certificate}
Write any combined kernel after cancellation as
$F(t)=\sum_{j,q}c_{j,q}(t+j)^{-q}$, so that
$\sum_jc_{j,q}=a$ for $q=s$ and zero otherwise. Using the Kronecker delta, define
\[
 \beta_{j,q}=\sum_{h=0}^j c_{h,q}-a\,\delta_{q,s},\qquad
 G(t)=\sum_{j=0}^{n-1}\sum_{q=1}^s\frac{\beta_{j,q}}{(t+j)^q}.
\]
Then the exact certificate is
\begin{equation}
 F(t)=\frac a{t^s}+G(t)-G(t+1),\qquad
 \sum_{k\ge1}F(k)=a\zeta(s)+G(1).
 \label{eq:telescoper}
\end{equation}
At the first pole the recovered coefficient is
$a\delta_{q,s}+\beta_{0,q}$; at an interior pole it is
$\beta_{j,q}-\beta_{j-1,q}$; at the last it is $-\beta_{n-1,q}$.
These are exactly the $c_{j,q}$. The rational difference has no poles
and vanishes at infinity, proving the first identity. Telescoping and
$G(k)\to0$ prove the second. This identifies the rational constant
constructively; it does not itself improve its arithmetic.

The actual block deformation also has the exact recurrence
\[
 T_{n,m+1}(t)=
 \frac{(t-m-1)(t+n+m+1)}{(m+1)^2}\,T_{n,m}(t).
\]
Lean verifies the unnormalized kernel version. These are identities in the
summation variable and block parameter, not a new recurrence in $n$.

\paragraph{Verified scope.}
The exact audit checks all $m\le n$ at weights $3,5,7,9$ and
$n=0,2,4,6,8,12,20$: $236$ kernels and $18\,600$ coefficient assertions.
It also certifies $19$ combined-form ledgers, all with one-sign weights,
including every pole coefficient of their telescopers and the equality
of $G(1)$ with the rational constant. Changing the clearing range leaves
their primitive costs exactly unchanged. Five new Lean foundation
reports use only standard foundations, with no added axioms or placeholders.
The general denominator theorem and telescoper derivation are paper proofs.

These finite certificates do not prove an asymptotic cost bound. For example,
the weight-five ledger with $m=n,n-1$ has
$\log K_n/n\approx7.286432$ at $n=20$, compared with the limiting analytic
decay rate $-\lambda_5\approx4.970794$. This is a finite comparison, not an
impossibility theorem. The winning inequality in \eqref{eq:cost} remains
unproved, so no all-odd irrationality conclusion follows.

\begin{thebibliography}{9}
\small
\bibitem{Companion}
\emph{Explicit symmetric Pad\'e kernels and exact odd-zeta cancellation},
companion research note, Section 5.
\url{https://github.com/ranbirsurii/irrationality-perhaps-/blob/74629b3/notes/symmetric-pade-note.pdf}.
\bibitem{Zudilin}
W. Zudilin, \emph{Arithmetic of linear forms involving odd zeta values},
Journal de Th\'eorie des Nombres de Bordeaux 16 (2004), 251--291.
\url{https://www.numdam.org/item/JTNB_2004__16_1_251_0.pdf}.
\end{thebibliography}
\end{document}
